% paper/sections/02b_csf.tex
%
% §2 支配方程式の続き（02_governing.tex からの分割）
% 内容：表面張力の離散化 — Ghost Fluid Method (GFM) 定式化
%       CSF モデルの完全な定式化は付録~\ref{app:csf_full} を参照．
%
% ── 後方互換ラベル（他章からの参照を維持するため） ──────────────────────────────
% 本節以前に定義されていた方程式ラベルを保持する．
% 完全な定式化は付録~\ref{app:csf_full} に移動した．
\phantomsection\label{eq:young-laplace}%
\phantomsection\label{eq:Heps_def_preview}%
\phantomsection\label{eq:delta_eps_def}%
% ═══════════════════════════════════════════════════════════════════════════════
\subsection{表面張力の離散化}
\label{sec:csf}
% ═══════════════════════════════════════════════════════════════════════════════

表面張力は界面 $\Gamma$ 上の面力であり，NS 方程式（体積方程式）への組み込みには
何らかの離散化方針が必要である．
標準的な手法は Brackbill \textit{et al.}~\cite{Brackbill1992} による
\textbf{連続表面力（CSF: Continuum Surface Force）モデル}であり，
表面張力を幅 $O(\varepsilon)$ の平滑化領域に分散した体積力
$\bm{f}_\sigma = \sigma\kappa\bnabla\psi$ として扱う
（完全な定式化は付録~\ref{app:csf_full}）．
CSF モデルは $O(\varepsilon^2) \approx O(\Delta x^2)$ のモデル誤差を持ち，
本稿の空間精度の律速段階となる（§\ref{sec:pressure_accuracy_summary}）．

本稿では，この限界を根本的に除去するため，
界面を数学的不連続面として扱う \textbf{Ghost Fluid Method (GFM)} \cite{Fedkiw1999} を採用する．

% ─────────────────────────────────────────────────────────────────────────────
\subsubsection{Ghost Fluid Method (GFM) による鋭利な界面処理}
\label{sec:gfm}
% ─────────────────────────────────────────────────────────────────────────────

GFM \cite{Fedkiw1999} は界面 $\Gamma$ を数学的不連続面として扱い，表面張力を PPE の右辺に直接組み込む手法である．
CSF 体積力は予測子ステップ（Step~5）から完全に除去され，代わりにジャンプ条件が PPE の右辺ベクトルへの補正項として組み込まれる．
これにより，CSF モデル誤差 $\Ord{\varepsilon^2}$ が本質的に除去される．

\paragraph{ジャンプ条件}

界面 $\Gamma$ における圧力の一般ジャンプ条件は $[p] = \sigma\kappa + 2[\mu](\nabla\mathbf{u}\cdot\mathbf{n})\cdot\mathbf{n}$ であるが，
本稿では粘性ジャンプ $[\mu] \equiv \mu_1 - \mu_2 \neq 0$ による追加項は界面法線速度勾配と乗算されるため，
低速流れの代表的計算ではその寄与が $\sigma\kappa$ に比べて小さいと仮定し，単純化ジャンプ条件：
\begin{equation}
  [p]_\Gamma \equiv p_1 - p_2 = \frac{\kappa}{We}
  \label{eq:gfm_jump}
\end{equation}
を採用する（$p_1$: 液相側圧力，$p_2$: 気相側圧力；無次元形式）．
この近似の妥当性は $We$ および密度比に依存し，高密度比・高 Reynolds 数系では再検討が必要である．

\paragraph{PPE への組み込み（ゴースト圧力置換）}

GFM では，界面を挟む格子点ペア $(i,\,i+1)$ において，相 1 側（$\phi_i > 0$）の PPE ステンシルで参照する相 2 ノード $i+1$ の圧力を，ジャンプ条件を加えた「ゴースト値」で置換する：
\begin{equation}
  \tilde{p}_{i+1}^{\text{ghost}} = p_{i+1} + [p]_\Gamma = p_{i+1} + \frac{\kappa_f}{We}
  \label{eq:gfm_ghost_p}
\end{equation}
ここで $\kappa_f$ は界面における曲率（フェイス補間値）である．この置換により，PPE の行列係数は変化せず，右辺ベクトル $\mathbf{b}$ に補正項が加わる形となる：
\begin{equation}
  b_i^{\mathrm{GFM}} = b_i - \left(\frac{1}{\rho}\right)_{i+1/2}^{\!\mathrm{harm}} \frac{\kappa_f}{We \cdot h^2}
  \quad \text{（界面を右に持つセル $i$ に対して）}
  \label{eq:gfm_rhs_correction}
\end{equation}
ここで $(1/\rho)_{i+1/2}^{\mathrm{harm}} = 2/(\rho_i + \rho_{i+1})$ は界面フェイスの調和平均密度逆数（PPE 演算子との整合性；付録~\ref{app:rhie_chow_full} 参照）であり，符号は界面の向き（$\phi_i > 0$ か $\phi_i < 0$ か）に従い，各軸方向について独立に適用する．

\paragraph{速度場の数値拡張（Extension）}

予測子ステップで CSF 体積力が除去されるため，対応して $\bu^*$ の PPE 右辺は CCD によるセル中心発散
\begin{equation}
  q_h = \frac{1}{\Delta t}\left(D_x^{(1)} u^* + D_y^{(1)} v^*\right)
  \label{eq:gfm_ccd_div}
\end{equation}
を使用する（Rhie--Chow 面速度補間は不要；チェッカーボード抑制は §\ref{sec:dccd_decoupling} の DCCD 散逸機構による）．

\paragraph{CCD との数学的相性}

本節で示した GFM 定式化（式~\eqref{eq:gfm_ghost_p}--\eqref{eq:gfm_rhs_correction}）は，圧力のゴースト値置換のみを扱う．
速度場を界面越しに滑らかに拡張する \textbf{Extension PDE}（$\partial q/\partial\tau \pm \mathbf{n}\cdot\nabla q = 0$）を各相について解くことで，
CCD のステンシル $(i-2,\ldots,i+2)$ が界面直近まで $C^\infty$ に近い滑らかな場を参照できるようになり，
\textbf{界面直近まで $\Ord{h^6}$ 精度が維持される}可能性がある．
Extension PDE の実装は本稿の現スコープ外であり，今後の課題とする（§\ref{sec:future_work}）．

% ═══════════════════════════════════════════════════════════════════════════════
